\documentclass[12pt]{article}
\usepackage[utf8]{inputenc}
\usepackage[T1]{fontenc}
\usepackage{hyperref}
\title{Saleman Language Resources for K-12 Curriculum Development}
\author{Compiled by Claude Sonnet 5.5 (claude-sonnet-5-5) for Andr\'as Kornai}
\date{}
\begin{document}
\maketitle

\section{Introduction}
Saleman (ISO 639-3 sau) is an Austronesian language of north-western Seram, Central Maluku, Indonesia. It is spoken in Negeri Saleman, Seram Utara Barat. The speaker count is uncertain. Ethnologue gives about 4,800 from 1989 data, while UNESCO's atlas lists 1,671 speakers from BPS 2018 data. Very little documentation was found.

\section{Primary Educational Resources}
No primers, textbooks or school programmes for Saleman were found. UNESCO says the language has no established orthography.

\section{Linguistic and Reference Materials}
\begin{itemize}
\item Wikipedia entry on the Saleman language: \url{https://en.wikipedia.org/wiki/Saleman_language}
\item UNESCO World Atlas of Languages entry: \url{https://ar.wal.unesco.org/countries/indonesia/languages/saleman}
\item Mardi Nugroho, language vitality study of Saleman in Jurnal Ranah (2020). It finds vitality declining and documentation a weak factor: \url{https://ojs.badanbahasa.kemdikbud.go.id/jurnal/index.php/jurnal_ranah/article/download/2938/1496}
\end{itemize}

\section{Religious and Community Resources}
Joshua Project profile of the Saleman people. Its religion figures are internally inconsistent: \url{https://joshuaproject.net/people_groups/14661}

\section{Digital and Audio Resources}
None found.

\section{Curriculum Development Context}
Saleman is declining and under-documented, with no known orthography or teaching materials. Any curriculum would first need documentation work: wordlists, recordings and an agreed orthography, with the community involved. Indonesian would remain the medium of instruction in the meantime.
\end{document}